\documentclass[11pt]{article}

\usepackage[utf8]{inputenc}
\usepackage[T1]{fontenc}
\usepackage{amsmath}
\usepackage{amssymb}
\usepackage{booktabs}
\usepackage{geometry}
\usepackage{longtable}
\usepackage{array}
\usepackage{xcolor}
\usepackage{hyperref}

\geometry{margin=1in}
\hypersetup{
  colorlinks=true,
  linkcolor=blue!55!black,
  urlcolor=blue!55!black,
  citecolor=blue!55!black,
  pdftitle={Comprehensive Validation Analysis: Thesis Note},
  pdfauthor={TFM Bayesian Football Analytics Pipeline}
}

\setlength{\parindent}{0pt}
\setlength{\parskip}{0.55em}
\renewcommand{\arraystretch}{1.12}
\setlength{\tabcolsep}{3pt}
\emergencystretch=2em
\newcolumntype{L}[1]{>{\raggedright\arraybackslash}p{#1}}
\newcolumntype{R}[1]{>{\raggedleft\arraybackslash}p{#1}}
\newcommand{\code}[1]{\texttt{\detokenize{#1}}}
\newcommand{\dd}{\,\mathrm{d}}

\title{\textbf{Comprehensive Validation Analysis}\\
\large Thesis-Ready Interpretation of the Additive Bayesian Modelling Pipeline}
\author{Bayesian Football Analytics Pipeline}
\date{Generated 2026-05-14}

\begin{document}
\maketitle

\begin{abstract}
This note summarizes the validation state of the real diagonal, simulated
diagonal, and simulated low-rank recovery models. The main thesis result is a
descriptive variance decomposition of multivariate football player-season
outcomes into player, season, and residual components. The real-data diagonal
additive model is usable with mild convergence caveats. Simulation results
support the diagonal variance-decomposition machinery for player and residual
components and, more cautiously, for season components. The low-rank recovery
experiment is an important negative result: it samples without Hamiltonian
Monte Carlo geometry failures but does not identify loading matrices reliably.
PSIS-LOO was computed, but Pareto-\(k\) diagnostics make the approximation
unreliable for formal model selection.
\end{abstract}

\section{Motivation and Research Question}

The scientific aim is not simply to fit a model that runs. The aim is to
establish which football-performance quantities can be interpreted as stable
player-level variation, which can be interpreted as season-level variation, and
which remain mostly unexplained row-level variation after accounting for
players and seasons.

The central research question is:
\begin{quote}
For a multivariate panel of standardized football features, how much variation
is attributable to stable player differences, common season shifts, and
player-season residual variation?
\end{quote}

The analysis conditions on the observed Wyscout-derived player-season panel.
It does not model transfers, injuries, tactical systems, position, team
identity, coach identity, match-level opponent strength, or minutes selection
beyond the existing filtering. Therefore a player effect should be interpreted
as a stable player-season profile effect within the observed data-generating
and selection process, not as a pure causal ability parameter.

\section{Data Structure and Estimands}

\subsection{Panel Structure}

The modelling unit is an observed player-season row. Missing player-seasons are
not imputed. This choice avoids assuming that absence from the panel is
ignorable. In football data, absence may be caused by transfers, injuries,
relegation, age, retirement, or selection. The cost is that conclusions apply
to observed eligible player-seasons, not to all possible player-season
combinations.

\begin{table}[htbp]
\centering
\small
\caption{Observed model object summary.}
\label{tab:model-object}
\begin{tabular}{L{0.58\textwidth}R{0.22\textwidth}}
\toprule
Quantity & Value \\
\midrule
Observed player-season rows & 4,944 \\
Selected features & 53 \\
Players & 1,650 \\
Seasons & 12 \\
One-season players & 605 \\
Repeated players & 1,045 \\
Players with gaps & 346 \\
Feature matrix rank & 53 \\
Maximum absolute feature correlation & 0.972 \\
\bottomrule
\end{tabular}
\end{table}

The feature matrix is full rank. The maximum retained absolute feature
correlation is high, mainly reflecting closely related pass-volume variables.
This does not break the diagonal additive model, but it warns that some
features carry nearly redundant information.

\subsection{Notation}

\begin{longtable}{L{0.20\textwidth}L{0.72\textwidth}}
\caption{Core notation used in the validation analysis.}
\label{tab:notation}\\
\toprule
Symbol & Meaning \\
\midrule
\endfirsthead
\toprule
Symbol & Meaning \\
\midrule
\endhead
\bottomrule
\endfoot
\(n = 1,\ldots,N\) & Observed player-season row index. \\
\(p = 1,\ldots,P\) & Feature index. \\
\(i = 1,\ldots,I\) & Player index. \\
\(j = 1,\ldots,S\) & Season index. \\
\(N\) & Number of observed player-season rows; here \(N=4{,}944\). \\
\(P\) & Number of selected scaled football features; here \(P=53\). \\
\(I\) & Number of players; here \(I=1{,}650\). \\
\(S\) & Number of seasons; here \(S=12\). \\
\(i[n]\) & Player associated with row \(n\). \\
\(j[n]\) & Season associated with row \(n\). \\
\(\mathbf{y}_n\in\mathbb{R}^{P}\) & Full multivariate response for row \(n\). \\
\(y_{np}\) & Value of feature \(p\) in row \(n\). \\
\(\boldsymbol{\mu}\in\mathbb{R}^{P}\) & Population mean in scaled feature space. \\
\(\mathbf{a}_i\in\mathbb{R}^{P}\) & Player effect for player \(i\). \\
\(\mathbf{b}_j\in\mathbb{R}^{P}\) & Season effect for season \(j\). \\
\(\boldsymbol{\epsilon}_n\in\mathbb{R}^{P}\) & Residual player-season effect for row \(n\). \\
\(\boldsymbol{\Sigma}_a\) & Player-effect covariance matrix. \\
\(\boldsymbol{\Sigma}_b\) & Season-effect covariance matrix. \\
\(\boldsymbol{\Sigma}_{\epsilon}\) & Residual covariance matrix. \\
\(\boldsymbol{\Lambda}_a,\boldsymbol{\Lambda}_b\) & Player and season low-rank loading matrices. \\
\(\boldsymbol{\Psi}_a,\boldsymbol{\Psi}_b\) & Player and season uniqueness variance vectors. \\
\end{longtable}

The response vector is
\begin{equation}
  \mathbf{y}_n = (y_{n1},\ldots,y_{nP})^\top ,
\end{equation}
where each column has been standardized before modelling. Standardization makes
variance components comparable across heterogeneous football features such as
shots per 90 minutes, pass rates, recoveries, and card counts.

\subsection{Model Components}

\begin{longtable}{L{0.22\textwidth}L{0.68\textwidth}}
\caption{Model components and their interpretation.}
\label{tab:model-components}\\
\toprule
Component & Interpretation \\
\midrule
\endfirsthead
\toprule
Component & Interpretation \\
\midrule
\endhead
\bottomrule
\endfoot
\(\boldsymbol{\mu}\) & Population-level average scaled football profile. \\
\(\mathbf{a}_i\) & Stable player-specific deviation from the population profile. \\
\(\mathbf{b}_j\) & Common season-specific deviation shared across observed rows. \\
\(\boldsymbol{\epsilon}_n\) & Remaining row-level variation after player and season effects. \\
\(\sigma_{a,p}\) & Player-effect standard deviation for feature \(p\). \\
\(\sigma_{b,p}\) & Season-effect standard deviation for feature \(p\). \\
\(\sigma_{\epsilon,p}\) & Residual standard deviation for feature \(p\). \\
\(\operatorname{diag}(\boldsymbol{\sigma}^2)\) & Diagonal covariance matrix with one variance per feature. \\
\(\boldsymbol{\Lambda}\boldsymbol{\Lambda}^{\top}\) & Low-rank covariance contribution from latent factors. \\
\(\boldsymbol{\Psi}\) & Feature-specific uniqueness variance not explained by common factors. \\
\end{longtable}

\subsection{Estimands}

For each feature \(p\), the real diagonal additive model estimates
\begin{align}
  V_{\mathrm{player},p} &= \sigma_{a,p}^{2}, \\
  V_{\mathrm{season},p} &= \sigma_{b,p}^{2}, \\
  V_{\mathrm{residual},p} &= \sigma_{\epsilon,p}^{2}.
\end{align}

The proportional variance decomposition is
\begin{align}
  \pi_{\mathrm{player},p}
    &= \frac{V_{\mathrm{player},p}}
            {V_{\mathrm{player},p}+V_{\mathrm{season},p}+V_{\mathrm{residual},p}}, \\
  \pi_{\mathrm{season},p}
    &= \frac{V_{\mathrm{season},p}}
            {V_{\mathrm{player},p}+V_{\mathrm{season},p}+V_{\mathrm{residual},p}}, \\
  \pi_{\mathrm{residual},p}
    &= \frac{V_{\mathrm{residual},p}}
            {V_{\mathrm{player},p}+V_{\mathrm{season},p}+V_{\mathrm{residual},p}}.
\end{align}

These proportions are the most interpretable thesis outputs. They answer how
much of each feature's variation is stable across players, how much is shared
by seasons, and how much remains row-specific after accounting for player and
season. They are model-based variance components, not causal effects.

\section{Real Diagonal Additive Model}

The primary real-data model is
\begin{equation}
  \mathbf{y}_n =
  \boldsymbol{\mu} + \mathbf{a}_{i[n]} + \mathbf{b}_{j[n]} +
  \boldsymbol{\epsilon}_n .
  \label{eq:additive}
\end{equation}

The diagonal covariance assumptions are
\begin{align}
  \mathbf{a}_i &\sim \mathcal{N}\!\left(
    \mathbf{0},\operatorname{diag}(\boldsymbol{\sigma}_a^2)\right), \\
  \mathbf{b}_j &\sim \mathcal{N}\!\left(
    \mathbf{0},\operatorname{diag}(\boldsymbol{\sigma}_b^2)\right), \\
  \boldsymbol{\epsilon}_n &\sim \mathcal{N}\!\left(
    \mathbf{0},\operatorname{diag}(\boldsymbol{\sigma}_{\epsilon}^2)\right).
\end{align}

The player and season effects are non-centered and column-centered in the
pipeline so that \(\boldsymbol{\mu}\) remains identifiable. The model is a
deliberate first-order variance decomposition. Its strength is interpretability:
each feature receives a player, season, and residual variance share. It is not
intended to answer the separate question of latent cross-feature covariance
structure.

\section{Simulation Design}

The simulation uses the same row count, feature count, player mapping, and
season mapping as the observed panel.

\begin{table}[htbp]
\centering
\small
\caption{Simulation design summary.}
\label{tab:simulation-summary}
\begin{tabular}{L{0.58\textwidth}R{0.20\textwidth}}
\toprule
Quantity & Value \\
\midrule
Simulation seed & 20260913 \\
Rows & 4,944 \\
Features & 53 \\
Players & 1,650 \\
Seasons & 12 \\
Player low-rank dimension & 2 \\
Season low-rank dimension & 1 \\
\bottomrule
\end{tabular}
\end{table}

The simulated data are generated from known low-rank plus diagonal player and
season covariance components. This creates two validation exercises. First, the
diagonal additive model is fitted to knowingly misspecified low-rank data to
test whether diagonal variance components are still recovered. Second, a
low-rank model is fitted to evaluate whether the known loadings, uniquenesses,
and covariance diagonals can be recovered.

\section{Low-Rank Recovery Model}

The low-rank recovery model uses
\begin{align}
  \boldsymbol{\Sigma}_a
    &= \boldsymbol{\Lambda}_a\boldsymbol{\Lambda}_a^\top
       + \operatorname{diag}(\boldsymbol{\Psi}_a), \\
  \boldsymbol{\Sigma}_b
    &= \boldsymbol{\Lambda}_b\boldsymbol{\Lambda}_b^\top
       + \operatorname{diag}(\boldsymbol{\Psi}_b), \\
  \boldsymbol{\Sigma}_{\epsilon}
    &= \operatorname{diag}(\boldsymbol{\sigma}_{\epsilon}^{2}).
\end{align}

The model estimates player loadings \(\boldsymbol{\Lambda}_a\), season loadings
\(\boldsymbol{\Lambda}_b\), player uniqueness scales, season uniqueness scales,
residual scales, and the implied covariance diagonals. It uses pure-anchor
constraints: selected anchor variables have positive own-factor loadings and
zero off-factor loadings. These constraints reduce rotation and sign ambiguity,
but the results show that they do not fully solve practical identifiability.

\section{Diagnostics and Posterior Health}

\subsection{Diagnostic Definitions}

\begin{longtable}{L{0.23\textwidth}L{0.67\textwidth}}
\caption{Diagnostics and validation metrics.}
\label{tab:diagnostics-definitions}\\
\toprule
Diagnostic & Interpretation \\
\midrule
\endfirsthead
\toprule
Diagnostic & Interpretation \\
\midrule
\endhead
\bottomrule
\endfoot
\(\widehat{R}\) & Chain convergence diagnostic. Values close to 1 are good; values above 1.01 indicate possible mixing issues. \\
Bulk ESS & Effective sample size for the central part of the posterior. \\
Tail ESS & Effective sample size for posterior tails. \\
Divergence & HMC warning that trajectories had difficulty following posterior geometry. \\
Treedepth hit & NUTS reached the configured maximum tree depth. \\
E-BFMI & Energy Bayesian Fraction of Missing Information; low values indicate poor energy exploration. \\
PSIS-LOO & Pareto-smoothed importance-sampling approximation to leave-one-out cross-validation. \\
Pareto-\(k\) & Stability diagnostic for PSIS weights; \(k>0.7\) is problematic and \(k>1\) is severe. \\
MAE & Mean absolute error between posterior estimate and simulation truth. \\
RMSE & Root mean squared error between posterior estimate and simulation truth. \\
Coverage & Proportion of true values inside the central 90\% posterior interval. \\
\end{longtable}

Posterior diagnostics are part of the scientific argument. Parameters with
clean \(\widehat{R}\) and adequate ESS can be interpreted normally. Parameters
with mild warnings should be interpreted cautiously. Parameters with severe
warnings should not be used for substantive conclusions.

\subsection{Sampler Geometry}

\begin{table}[htbp]
\centering
\small
\caption{Sampler geometry summary for the three fitted models.}
\label{tab:sampler-geometry}
\begin{tabular}{L{0.23\textwidth}R{0.09\textwidth}R{0.08\textwidth}R{0.09\textwidth}R{0.10\textwidth}R{0.11\textwidth}R{0.11\textwidth}}
\toprule
Model & Draws & Div. & Max TD & TD hits & Min E-BFMI & Mean accept \\
\midrule
Real diagonal & 4,000 & 0 & 8 & 0 & 0.561 & 0.953 \\
Simulated diagonal & 4,000 & 0 & 8 & 0 & 0.551 & 0.951 \\
Simulated low-rank & 4,000 & 0 & 8 & 0 & 0.586 & 0.951 \\
\bottomrule
\end{tabular}
\end{table}

All three models have clean HMC geometry: no divergences, no treedepth
saturations at the configured limit of 12, E-BFMI above common warning
thresholds, and average acceptance rates close to the target
\(\texttt{adapt\_delta}=0.95\). The main problems are therefore not HMC
geometry problems. They are posterior mixing, identifiability, and predictive
validation diagnostics.

\subsection{Posterior Health}

\begin{table}[htbp]
\centering
\small
\caption{Posterior summary diagnostics.}
\label{tab:posterior-health}
\begin{tabular}{L{0.21\textwidth}R{0.12\textwidth}R{0.09\textwidth}R{0.08\textwidth}R{0.08\textwidth}R{0.11\textwidth}R{0.11\textwidth}}
\toprule
Model & Parameters & Max \(\widehat{R}\) & \(>1.01\) & \(>1.05\) & Min bulk ESS & Min tail ESS \\
\midrule
Real diagonal & 531 & 1.035 & 13 & 0 & 159 & 175 \\
Simulated diagonal & 531 & 1.010 & 3 & 0 & 779 & 1,412 \\
Simulated low-rank & 688 & 1.739 & 233 & 208 & 6 & 107 \\
\bottomrule
\end{tabular}
\end{table}

The real diagonal model is mostly usable, with mild mixing warnings in a small
number of player variance parameters. The simulated diagonal model is healthy.
The simulated low-rank model is not healthy for loading interpretation.

\begin{table}[htbp]
\centering
\small
\caption{Worst real diagonal posterior diagnostics.}
\label{tab:real-worst-diagnostics}
\begin{tabular}{L{0.17\textwidth}L{0.28\textwidth}R{0.08\textwidth}R{0.07\textwidth}R{0.08\textwidth}R{0.09\textwidth}R{0.09\textwidth}}
\toprule
Parameter & Feature & Mean & SD & \(\widehat{R}\) & Bulk ESS & Tail ESS \\
\midrule
\code{var_a[11]} & \code{per90_forward_passes} & 0.827 & 0.032 & 1.035 & 159 & 290 \\
\code{sigma_a[11]} & \code{per90_forward_passes} & 0.909 & 0.018 & 1.035 & 159 & 290 \\
\code{prop_e[11]} & \code{per90_forward_passes} & 0.150 & 0.006 & 1.028 & 202 & 642 \\
\code{prop_a[11]} & \code{per90_forward_passes} & 0.846 & 0.006 & 1.027 & 179 & 376 \\
\code{var_a[44]} & \code{per90_sliding_tackles} & 0.652 & 0.029 & 1.016 & 894 & 1,643 \\
\code{var_a[16]} & \code{per90_progressive_passes} & 0.752 & 0.029 & 1.015 & 509 & 1,746 \\
\code{var_a[24]} & \code{rate_successful_crosses} & 0.022 & 0.013 & 1.012 & 197 & 177 \\
\bottomrule
\end{tabular}
\end{table}

No real-data \(\widehat{R}\) exceeds 1.05. For thesis reporting these
parameters should be described as mildly under-mixed rather than perfectly
converged.

\begin{table}[htbp]
\centering
\small
\caption{Worst low-rank posterior diagnostics.}
\label{tab:lowrank-worst-diagnostics}
\begin{tabular}{L{0.43\textwidth}R{0.10\textwidth}R{0.10\textwidth}R{0.10\textwidth}R{0.12\textwidth}}
\toprule
Parameter & Mean & SD & \(\widehat{R}\) & Bulk ESS \\
\midrule
\code{lambda_b_free[37]} / \(\Lambda_b[38,1]\) & -0.010 & 0.271 & 1.739 & 6.1 \\
\code{lambda_a_free[41]} / \(\Lambda_a[43,1]\) & 0.000 & 0.111 & 1.738 & 6.2 \\
\code{lambda_a_free[23]} / \(\Lambda_a[25,1]\) & 0.000 & 0.210 & 1.738 & 6.1 \\
\code{lambda_b_free[24]} / \(\Lambda_b[25,1]\) & 0.018 & 0.328 & 1.738 & 6.1 \\
\code{psi_b[1]} & 0.332 & 0.246 & 1.737 & 6.2 \\
\code{psi_a[1]} & 0.418 & 0.127 & 1.737 & 6.2 \\
\bottomrule
\end{tabular}
\end{table}

These values are a strong identifiability and mixing warning. Chains are not
exploring the same posterior region for many loading and uniqueness parameters,
even though HMC trajectories themselves are clean.

\section{Real-Data Variance Decomposition}

The real diagonal additive model decomposes scaled feature variance into
player, season, and residual components. Averaged across features, the player
component is dominant.

\begin{table}[htbp]
\centering
\small
\caption{Mean variance shares across the 53 real-data features.}
\label{tab:mean-variance-shares}
\begin{tabular}{L{0.42\textwidth}R{0.18\textwidth}}
\toprule
Component & Mean share \\
\midrule
Player & 0.634 \\
Season & 0.047 \\
Residual & 0.319 \\
\bottomrule
\end{tabular}
\end{table}

This means that stable player-level heterogeneity is the dominant source of
variation in most selected features. Season-level variation exists but is much
smaller on average. Residual variation remains important, especially for rate
features.

\subsection{Features Dominated by Player Effects}

\begin{table}[htbp]
\centering
\small
\caption{Highest player-level variance shares.}
\label{tab:player-dominated}
\begin{tabular}{L{0.42\textwidth}R{0.13\textwidth}R{0.13\textwidth}R{0.13\textwidth}}
\toprule
Feature & Player & Season & Residual \\
\midrule
\code{per90_recoveries} & 0.884 & 0.016 & 0.100 \\
\code{per90_touch_in_box} & 0.872 & 0.001 & 0.127 \\
\code{per90_offensive_duels} & 0.869 & 0.056 & 0.075 \\
\code{per90_shots} & 0.860 & 0.008 & 0.132 \\
\code{per90_forward_passes} & 0.846 & 0.004 & 0.150 \\
\code{per90_crosses} & 0.846 & 0.006 & 0.148 \\
\code{per90_long_passes} & 0.843 & 0.005 & 0.152 \\
\code{per90_clearances} & 0.837 & 0.034 & 0.129 \\
\code{per90_passes} & 0.829 & 0.011 & 0.160 \\
\code{per90_back_passes} & 0.823 & 0.008 & 0.169 \\
\bottomrule
\end{tabular}
\end{table}

Volume and style variables are strongly player-specific. This is consistent
with stable player roles, positions, and tactical usage.

\subsection{Features With the Largest Season Effects}

\begin{table}[htbp]
\centering
\small
\caption{Highest season-level variance shares.}
\label{tab:season-dominated}
\begin{tabular}{L{0.42\textwidth}R{0.13\textwidth}R{0.13\textwidth}R{0.13\textwidth}}
\toprule
Feature & Player & Season & Residual \\
\midrule
\code{per90_smart_passes} & 0.440 & 0.376 & 0.183 \\
\code{per90_pressing_duels} & 0.455 & 0.320 & 0.225 \\
\code{per90_accelerations} & 0.598 & 0.164 & 0.238 \\
\code{per90_dribbles_against} & 0.693 & 0.136 & 0.171 \\
\code{rate_defensive_duels_won} & 0.244 & 0.129 & 0.627 \\
\code{per90_through_passes} & 0.583 & 0.121 & 0.296 \\
\code{per90_sliding_tackles} & 0.606 & 0.119 & 0.275 \\
\code{per90_loose_ball_duels} & 0.718 & 0.113 & 0.168 \\
\code{per90_dribbles} & 0.800 & 0.101 & 0.100 \\
\code{rate_successful_progressive_passes} & 0.484 & 0.084 & 0.432 \\
\bottomrule
\end{tabular}
\end{table}

The strongest season effects occur in features that plausibly reflect changing
tactical environments, league-wide style, event definitions, or tactical
fashion over seasons. The most prominent examples are smart passes and pressing
duels.

\subsection{Features Dominated by Residual Effects}

\begin{table}[htbp]
\centering
\small
\caption{Highest residual variance shares.}
\label{tab:residual-dominated}
\begin{tabular}{L{0.42\textwidth}R{0.13\textwidth}R{0.13\textwidth}R{0.13\textwidth}}
\toprule
Feature & Player & Season & Residual \\
\midrule
\code{rate_successful_crosses} & 0.022 & 0.003 & 0.975 \\
\code{rate_successful_sliding_tackles} & 0.057 & 0.011 & 0.932 \\
\code{rate_shots_on_target} & 0.167 & 0.001 & 0.832 \\
\code{rate_goal_conversion} & 0.180 & 0.003 & 0.818 \\
\code{rate_dribbles_against_won} & 0.184 & 0.008 & 0.808 \\
\code{rate_successful_dribbles} & 0.231 & 0.009 & 0.760 \\
\code{rate_defensive_duels_won} & 0.244 & 0.129 & 0.627 \\
\code{per90_assists} & 0.391 & 0.010 & 0.599 \\
\code{rate_offensive_duels_won} & 0.412 & 0.016 & 0.573 \\
\code{per90_yellow_cards} & 0.441 & 0.017 & 0.542 \\
\bottomrule
\end{tabular}
\end{table}

Rate variables and sparse outcome variables are noisy. They depend strongly on
opportunity counts and row-level context and are much less stable than
volume/style features.

\section{Simulation Recovery Results}

\subsection{Diagonal Model on Low-Rank Data}

The simulated diagonal model is intentionally misspecified because the data
were generated from low-rank plus diagonal covariance. Despite this, it recovers
diagonal variance components well.

\begin{table}[htbp]
\centering
\small
\caption{Recovery metrics for the simulated diagonal additive model.}
\label{tab:sim-diagonal-recovery}
\begin{tabular}{L{0.20\textwidth}R{0.12\textwidth}R{0.12\textwidth}R{0.12\textwidth}R{0.14\textwidth}R{0.14\textwidth}}
\toprule
Component & Bias & MAE & RMSE & Corr. & 90\% coverage \\
\midrule
Player & 0.0028 & 0.0078 & 0.0112 & 0.996 & 0.906 \\
Residual & -0.0009 & 0.0040 & 0.0055 & 0.998 & 0.925 \\
Season & 0.0250 & 0.0300 & 0.0468 & 0.885 & 0.868 \\
\bottomrule
\end{tabular}
\end{table}

Player variance recovery is excellent, residual variance recovery is excellent,
and season variance recovery is weaker but still useful. The season component
is hardest because there are only 12 seasons, which is a small number of
season-level units for feature-by-feature season variance estimation.

\begin{table}[htbp]
\centering
\small
\caption{Worst season variance errors in the simulated diagonal model.}
\label{tab:sim-diagonal-worst-season}
\begin{tabular}{L{0.48\textwidth}R{0.16\textwidth}R{0.14\textwidth}R{0.12\textwidth}}
\toprule
Feature & Posterior mean & Truth & Error \\
\midrule
\code{rate_successful_sliding_tackles} & 0.354 & 0.139 & 0.214 \\
\code{per90_xg_shot} & 0.221 & 0.103 & 0.118 \\
\code{per90_goals} & 0.236 & 0.151 & 0.085 \\
\code{per90_fouls} & 0.144 & 0.064 & 0.080 \\
\code{per90_defensive_duels} & 0.130 & 0.057 & 0.073 \\
\bottomrule
\end{tabular}
\end{table}

The bias is mostly upward for season variance, which is consistent with a
diagonal model approximating low-rank seasonal covariance using feature-specific
season variance terms.

\subsection{Low-Rank Loading Recovery}

\begin{table}[htbp]
\centering
\small
\caption{Low-rank loading recovery metrics.}
\label{tab:lowrank-loading-recovery}
\begin{tabular}{L{0.22\textwidth}R{0.14\textwidth}R{0.14\textwidth}R{0.16\textwidth}R{0.16\textwidth}}
\toprule
Parameter & MAE & RMSE & Corr. & 90\% coverage \\
\midrule
\(\Lambda_a\) & 0.082 & 0.131 & 0.804 & 0.840 \\
\(\Lambda_b\) & 0.102 & 0.124 & 0.262 & 1.000 \\
\bottomrule
\end{tabular}
\end{table}

The player loading matrix has moderate correlation with truth but large errors
and poor mixing. The season loading matrix is not recovered well: its
correlation with truth is only 0.262, and coverage is high because intervals
are too wide, not because estimates are precise.

\begin{table}[htbp]
\centering
\small
\caption{Worst \(\Lambda_a\) loading errors.}
\label{tab:lambda-a-worst}
\begin{tabular}{L{0.39\textwidth}R{0.08\textwidth}R{0.10\textwidth}R{0.10\textwidth}R{0.10\textwidth}R{0.08\textwidth}}
\toprule
Feature & Factor & Mean & Truth & Error & \(\widehat{R}\) \\
\midrule
\code{per90_shot_assists} & 1 & -0.003 & 0.491 & -0.494 & 1.736 \\
\code{per90_fouls} & 1 & -0.001 & 0.352 & -0.352 & 1.736 \\
\code{per90_shots_blocked} & 1 & -0.001 & 0.340 & -0.342 & 1.735 \\
\code{per90_opponent_half_recoveries} & 1 & 0.002 & -0.335 & 0.337 & 1.734 \\
\code{per90_long_passes} & 1 & 0.002 & -0.289 & 0.291 & 1.735 \\
\bottomrule
\end{tabular}
\end{table}

\begin{table}[htbp]
\centering
\small
\caption{Worst \(\Lambda_b\) loading errors.}
\label{tab:lambda-b-worst}
\begin{tabular}{L{0.39\textwidth}R{0.08\textwidth}R{0.10\textwidth}R{0.10\textwidth}R{0.10\textwidth}R{0.08\textwidth}}
\toprule
Feature & Factor & Mean & Truth & Error & \(\widehat{R}\) \\
\midrule
\code{rate_successful_sliding_tackles} & 1 & -0.042 & 0.298 & -0.340 & 1.736 \\
\code{per90_xg_shot} & 1 & -0.008 & 0.230 & -0.238 & 1.737 \\
\code{per90_passes} & 1 & -0.003 & 0.230 & -0.233 & 1.733 \\
\code{per90_dangerous_own_half_losses} & 1 & -0.008 & 0.206 & -0.214 & 1.735 \\
\code{per90_through_passes} & 1 & 0.023 & -0.185 & 0.208 & 1.734 \\
\bottomrule
\end{tabular}
\end{table}

Many free loadings have posterior means near zero, wide intervals, and very
high \(\widehat{R}\). The current low-rank model should therefore be treated as
a failed or incomplete loading-recovery experiment, not as a substantive
football-factor model.

\subsection{Low-Rank Variance and Scale Recovery}

\begin{table}[htbp]
\centering
\small
\caption{Low-rank variance and scale recovery metrics.}
\label{tab:lowrank-variance-recovery}
\begin{tabular}{L{0.26\textwidth}R{0.13\textwidth}R{0.13\textwidth}R{0.15\textwidth}R{0.15\textwidth}}
\toprule
Parameter & MAE & RMSE & Corr. & 90\% coverage \\
\midrule
\(\operatorname{diag}(\Sigma_a)\) & 0.009 & 0.013 & 0.997 & 0.811 \\
\(\operatorname{diag}(\Sigma_b)\) & 0.057 & 0.093 & 0.886 & 0.585 \\
\(\psi_a\) SD & 0.011 & 0.022 & 0.941 & 0.925 \\
\(\psi_b\) SD & 0.040 & 0.052 & 0.740 & 0.925 \\
\(\sigma_{\epsilon}\) SD & 0.004 & 0.005 & 0.998 & 0.925 \\
\bottomrule
\end{tabular}
\end{table}

The low-rank model recovers residual scales very well, player covariance
diagonals very well, and player uniqueness scales reasonably well. It does not
recover season covariance diagonals precisely, season uniqueness scales as well
as player scales, or loading matrices reliably.

\begin{table}[htbp]
\centering
\small
\caption{Worst \(\operatorname{diag}(\Sigma_b)\) errors in the low-rank recovery model.}
\label{tab:sigma-b-worst}
\begin{tabular}{L{0.43\textwidth}R{0.11\textwidth}R{0.10\textwidth}R{0.10\textwidth}R{0.09\textwidth}}
\toprule
Feature & Mean & Truth & Error & \(\widehat{R}\) \\
\midrule
\code{rate_successful_sliding_tackles} & 0.572 & 0.139 & 0.432 & 1.051 \\
\code{per90_goals} & 0.435 & 0.151 & 0.284 & 1.029 \\
\code{per90_xg_shot} & 0.342 & 0.103 & 0.238 & 1.009 \\
\code{per90_through_passes} & 0.219 & 0.084 & 0.134 & 1.013 \\
\code{per90_fouls} & 0.160 & 0.064 & 0.096 & 1.003 \\
\bottomrule
\end{tabular}
\end{table}

Recovering a covariance diagonal is easier than recovering a loading matrix.
Different loading configurations can imply similar total feature variance, so
good recovery of \(\operatorname{diag}(\Sigma_a)\) does not imply good recovery
of \(\Lambda_a\).

\section{PSIS-LOO Interpretation}

PSIS-LOO was computed with \code{BFA_RUN_LOO=true}. The nominal estimates are
shown in Table~\ref{tab:loo}. The simulated low-rank model has better nominal
predictive fit than the simulated diagonal model, with
\(\Delta\mathrm{elpd}_{\mathrm{low-rank}-\mathrm{diagonal}}\approx 5{,}576\).
This direction is expected because the simulated data were generated from a
low-rank process.

\begin{table}[htbp]
\centering
\small
\caption{PSIS-LOO results and Pareto-\(k\) diagnostics.}
\label{tab:loo}
\begin{tabular}{L{0.18\textwidth}R{0.11\textwidth}R{0.08\textwidth}R{0.10\textwidth}R{0.11\textwidth}R{0.08\textwidth}R{0.08\textwidth}R{0.08\textwidth}}
\toprule
Model & elpd\(_{\mathrm{loo}}\) & SE & \(p_{\mathrm{loo}}\) & looic & \(k>0.7\) & \(k>1\) & Max \(k\) \\
\midrule
Real diagonal & -224,673.4 & 1,007.0 & 49,545.7 & 449,346.9 & 4,354 & 3,088 & 3.82 \\
Simulated diagonal & -201,763.1 & 394.7 & 41,871.8 & 403,526.1 & 4,529 & 2,878 & 2.88 \\
Simulated low-rank & -196,186.8 & 363.3 & 32,963.7 & 392,373.7 & 4,215 & 2,104 & 2.08 \\
\bottomrule
\end{tabular}
\end{table}

The Pareto-\(k\) diagnostics are unacceptable: thousands of rows have
\(k>0.7\), thousands have \(k>1\), and maximum \(k\) is far above 1 for all
models. The LOO estimates should therefore not be used as formal model-selection
evidence. They are better interpreted as evidence that many rows are highly
influential under conditional leave-one-row-out scoring.

This is unsurprising for a hierarchical player-season model. Many players
appear only once, player effects can be informed heavily by a single row, and
leaving out that row can substantially change the effective posterior. Exact
\(K\)-fold cross-validation is preferable for real predictive validation.

\section{Limitations}

The main limitations are as follows.

\begin{enumerate}
\item The real diagonal model has mild mixing issues. The maximum
\(\widehat{R}\) is 1.035, 13 summarized parameters exceed 1.01, the minimum
bulk ESS is 159, and the minimum tail ESS is 175.
\item The low-rank recovery model has severe loading-level non-identification.
The maximum \(\widehat{R}\) is 1.739, many loadings have bulk ESS near 6, and
many loading posterior means are near zero despite nonzero truth.
\item Season effects are hard to estimate. There are only 12 seasons, and the
simulation confirms weaker recovery for season variance and especially season
loadings.
\item PSIS-LOO fails Pareto-\(k\) diagnostics and should not be used as decisive
model-selection evidence.
\item Memory and output size constrain diagnostics. Full fit CSVs are large,
and full draw extraction can hit R memory limits.
\end{enumerate}

\section{Proposed Fixes and Future Experiments}

Several extensions follow directly from the validation evidence.

\begin{enumerate}
\item Run a longer real diagonal model with more sampling iterations to check
whether the few weak variance parameters improve.
\item Build exact \(K\)-fold validation. Leave-player-out folds target
generalization to unseen players; leave-season-out folds target generalization
to new seasons; random row folds are appropriate only if the target is another
observed row from the same player-season process.
\item For simulation only, initialize low-rank parameters near truth to
distinguish poor exploration from deeper non-identifiability.
\item Strengthen low-rank anchor constraints by lower-bounding anchor loadings,
using stronger anchor priors, selecting higher-signal anchors, and adding
stronger shrinkage to free loadings.
\item Simplify the season structure. With only 12 seasons, diagonal season
effects or a strongly pooled scalar season score may be more defensible than a
feature-level low-rank season loading vector.
\item Add feature sensitivity checks: remove highly correlated pass-volume
duplicates, analyze rate-only and volume-only feature sets separately, stratify
by position if metadata are available, and test higher minimum-minute
thresholds.
\item Add targeted memory-safe diagnostics that extract only a small variable
set when trace plots or draw-level checks are needed.
\end{enumerate}

\section{Thesis-Ready Conclusions}

The pipeline has succeeded for descriptive additive variance decomposition.
The real diagonal additive model gives a defensible main result: stable player
effects dominate many volume and style features, season effects are generally
smaller but meaningful for selected tactical variables, and residual variation
dominates sparse rates and conversion or success variables.

The simulation validates the basic additive machinery. The diagonal model
recovers player and residual variance components very well, and season variance
moderately well, even when the simulated truth contains low-rank covariance.

The current low-rank recovery model should be reported as a diagnostic negative
result. It has clean sampler geometry but fails to identify loadings reliably,
especially for season effects. Its covariance-diagonal recovery is more
credible than its loading recovery, which reinforces the distinction between
variance-decomposition claims and latent-factor claims.

PSIS-LOO was computed but fails Pareto-\(k\) diagnostics. It should be
mentioned only as an attempted predictive diagnostic, not as decisive
model-selection evidence. Exact \(K\)-fold validation is the appropriate next
step for predictive comparison.

\end{document}
